\documentclass[10pt,a4paper]{amsart}
\usepackage{amsmath,amssymb,amsthm,graphicx}
% Retained preamble for the Stacks Project chapter snapshot in this directory.
% The chapter's own preamble.tex is a copy of the shared Stacks Project preamble; the block
% below is sliced MECHANICALLY, line for line, from it. Nothing was added, removed or
% redefined; the wrapper class, the xr-hyper/hyperref wiring and the externaldocument list
% are omitted because the assembler supplies the document class.
%
% lines taken from preamble.tex: 3, 5, 9, 11, 12, 156, 157, 158, ... (51 lines total)
%
% package coverage restored from the chapter's own preamble (2026-10-13): xy, xy-2cell

\usepackage[all]{xy}
\xyoption{2cell}
\UseAllTwocells

\newenvironment{reference}{\comment}{\endcomment}
\newenvironment{slogan}{\comment}{\endcomment}
\newenvironment{history}{\comment}{\endcomment}
\theoremstyle{plain}
\newtheorem{theorem}[subsection]{Theorem}
\newtheorem{proposition}[subsection]{Proposition}
\newtheorem{lemma}[subsection]{Lemma}
\theoremstyle{definition}
\newtheorem{definition}[subsection]{Definition}
\newtheorem{example}[subsection]{Example}
\newtheorem{exercise}[subsection]{Exercise}
\newtheorem{situation}[subsection]{Situation}
\theoremstyle{remark}
\newtheorem{remark}[subsection]{Remark}
\newtheorem{remarks}[subsection]{Remarks}
\numberwithin{equation}{subsection}
\def\lim{\mathop{\mathrm{lim}}\nolimits}
\def\colim{\mathop{\mathrm{colim}}\nolimits}
\def\Spec{\mathop{\mathrm{Spec}}}
\def\Hom{\mathop{\mathrm{Hom}}\nolimits}
\def\Ext{\mathop{\mathrm{Ext}}\nolimits}
\def\SheafHom{\mathop{\mathcal{H}\!\mathit{om}}\nolimits}
\def\SheafExt{\mathop{\mathcal{E}\!\mathit{xt}}\nolimits}
\def\Sch{\mathit{Sch}}
\def\Mor{\mathop{\mathrm{Mor}}\nolimits}
\def\Ob{\mathop{\mathrm{Ob}}\nolimits}
\def\Sh{\mathop{\mathit{Sh}}\nolimits}
\def\NL{\mathop{N\!L}\nolimits}
\def\CH{\mathop{\mathrm{CH}}\nolimits}
\def\proetale{{pro\text{-}\acute{e}tale}}
\def\etale{{\acute{e}tale}}
\def\QCoh{\mathit{QCoh}}
\def\Ker{\mathop{\mathrm{Ker}}}
\def\Im{\mathop{\mathrm{Im}}}
\def\Coker{\mathop{\mathrm{Coker}}}
\def\Coim{\mathop{\mathrm{Coim}}}
\DeclareMathSymbol{\boxtimes}{\mathbin}{AMSa}{"02}
\def\QCohstack{\mathcal{QC}\!\mathit{oh}}
\def\Cohstack{\mathcal{C}\!\mathit{oh}}
\def\Spacesstack{\mathcal{S}\!\mathit{paces}}
\def\Quotfunctor{\mathrm{Quot}}
\def\Hilbfunctor{\mathrm{Hilb}}
\def\Curvesstack{\mathcal{C}\!\mathit{urves}}
\def\Polarizedstack{\mathcal{P}\!\mathit{olarized}}
\def\Complexesstack{\mathcal{C}\!\mathit{omplexes}}
\def\Pic{\mathop{\mathrm{Pic}}\nolimits}
\def\Picardstack{\mathcal{P}\!\mathit{ic}}
\def\Picardfunctor{\mathrm{Pic}}
\def\Deformationcategory{\mathcal{D}\!\mathit{ef}}
\usepackage{cleveref}
\crefname{originallemma}{original lemma}{original lemmas}
\Crefname{originallemma}{Original lemma}{Original lemmas}
\makeatletter
\newlabel{algebra-lemma-complete-by-sub}{{090T}{source}{}{sourcepointer-1}{}}
\newlabel{algebra-lemma-complete-by-sub@cref}{{[originallemma][0][]090T}{[1][0][]source}{}{}{}}
\makeatother
\crefname{originallemma}{original lemma}{original lemmas}
\Crefname{originallemma}{Original lemma}{Original lemmas}
\makeatletter
\newlabel{formal-spaces-lemma-map-into-algebraic-space}{{0AQH}{source}{}{sourcepointer-2}{}}
\newlabel{formal-spaces-lemma-map-into-algebraic-space@cref}{{[originallemma][0][]0AQH}{[1][0][]source}{}{}{}}
\makeatother
\crefname{originallemma}{original lemma}{original lemmas}
\Crefname{originallemma}{Original lemma}{Original lemmas}
\makeatletter
\newlabel{lemma-effective}{{07X8}{source}{}{sourcepointer-3}{}}
\newlabel{lemma-effective@cref}{{[originallemma][0][]07X8}{[1][0][]source}{}{}{}}
\makeatother
\crefname{originallemma}{original lemma}{original lemmas}
\Crefname{originallemma}{Original lemma}{Original lemmas}
\makeatletter
\newlabel{lemma-rs}{{0GI6}{source}{}{sourcepointer-4}{}}
\newlabel{lemma-rs@cref}{{[originallemma][0][]0GI6}{[1][0][]source}{}{}{}}
\makeatother
\crefname{originallemma}{original lemma}{original lemmas}
\Crefname{originallemma}{Original lemma}{Original lemmas}
\makeatletter
\newlabel{restricted-lemma-base-change-formal-modification}{{0GDN}{source}{}{sourcepointer-5}{}}
\newlabel{restricted-lemma-base-change-formal-modification@cref}{{[originallemma][0][]0GDN}{[1][0][]source}{}{}{}}
\makeatother
\crefname{originallemma}{original lemma}{original lemmas}
\Crefname{originallemma}{Original lemma}{Original lemmas}
\makeatletter
\newlabel{restricted-theorem-dilatations}{{0GDU}{source}{}{sourcepointer-6}{}}
\newlabel{restricted-theorem-dilatations@cref}{{[originallemma][0][]0GDU}{[1][0][]source}{}{}{}}
\makeatother
\crefname{originallemma}{original lemma}{original lemmas}
\Crefname{originallemma}{Original lemma}{Original lemmas}
\makeatletter
\newlabel{situation-contractions}{{0GH8}{source}{}{sourcepointer-7}{}}
\newlabel{situation-contractions@cref}{{[originallemma][0][]0GH8}{[1][0][]source}{}{}{}}
\makeatother
\crefname{originallemma}{original lemma}{original lemmas}
\Crefname{originallemma}{Original lemma}{Original lemmas}
\makeatletter
\newlabel{spaces-more-morphisms-lemma-algebraize-morphism}{{0A4Z}{source}{}{sourcepointer-8}{}}
\newlabel{spaces-more-morphisms-lemma-algebraize-morphism@cref}{{[originallemma][0][]0A4Z}{[1][0][]source}{}{}{}}
\makeatother
\title{Effectiveness of formal objects when X' -> S is separated}
\author{The Stacks Project Authors}
\date{stacks-project:artin.tex:local-snapshot-20261012; GNU FDL 1.2 or later}
\begin{document}
\maketitle
\noindent\textit{Editorial note.} Statement span is lines 6236-6240 of sources/stacks-artin/pinned-original-source.tex; proof span is lines 6242-6327, the complete author proof (84 body lines).\par\medskip
\noindent\textit{Editorial note.} Navigation stubs: 8 cited-but-not-included labels are supplied as non-load-bearing pointers with their pinned tags where the snapshot tag index has them.\par\medskip
\noindent\textit{Editorial note.} extra\_packages is empty on purpose: the retained preamble already loads xy with options and a second declaration clashes in the assembler wrapper.\par\medskip
\noindent\textit{Editorial note.} Text is retained verbatim, including the author's own typos and placeholders.\par\medskip
\noindent\textit{Editorial note.} Queue index 10 of queue\_artin.json (pb=3285, stmt 6236-6240, proof 6242-6327).\par\medskip
\noindent\textit{Editorial note.} stacks-artin wave ar01 (2026-10-14): fresh adjudication accept per lane-s3/artin\_judgments.md; pool per Chief 10:5x direct dispatch (idle-first over S12's 01:5x note, SMC precedent); artin.tex 99 rows/67 with proof judged 46 accept / 21 reject / 5 AGP-overlap avoidance rows; patch-sd-envs contract spans scanned.\par\medskip
\noindent\textit{Allowed standard prerequisites (editorial).}\begin{itemize}
\item Explicit prerequisite boundary. The cited inputs are prerequisites; the argument in the retained span is not replaced by any cited result.
\end{itemize}\medskip
\ifdefined\hypertarget\hypertarget{sourcepointer-1}{}\fi
\noindent\textit{Source reference.} Original lemma 090T: sources/stacks-artin/pinned-original-source.tex, source lines 6327--6327. Cited by the retained statement or proof but not part of the collected spans; supplied only as a navigation pointer. The argument does not depend on its content.\par\medskip
\ifdefined\hypertarget\hypertarget{sourcepointer-2}{}\fi
\noindent\textit{Source reference.} Original lemma 0AQH: sources/stacks-artin/pinned-original-source.tex, source lines 6327--6327. Cited by the retained statement or proof but not part of the collected spans; supplied only as a navigation pointer. The argument does not depend on its content.\par\medskip
\ifdefined\hypertarget\hypertarget{sourcepointer-3}{}\fi
\noindent\textit{Source reference.} Original lemma 07X8: sources/stacks-artin/pinned-original-source.tex, source lines 1323--1326. Cited by the retained statement or proof but not part of the collected spans; supplied only as a navigation pointer. The argument does not depend on its content.\par\medskip
\ifdefined\hypertarget\hypertarget{sourcepointer-4}{}\fi
\noindent\textit{Source reference.} Original lemma 0GI6: sources/stacks-artin/pinned-original-source.tex, source lines 6175--6178. Cited by the retained statement or proof but not part of the collected spans; supplied only as a navigation pointer. The argument does not depend on its content.\par\medskip
\ifdefined\hypertarget\hypertarget{sourcepointer-5}{}\fi
\noindent\textit{Source reference.} Original lemma 0GDN: sources/stacks-artin/pinned-original-source.tex, source lines 6327--6327. Cited by the retained statement or proof but not part of the collected spans; supplied only as a navigation pointer. The argument does not depend on its content.\par\medskip
\ifdefined\hypertarget\hypertarget{sourcepointer-6}{}\fi
\noindent\textit{Source reference.} Original lemma 0GDU: sources/stacks-artin/pinned-original-source.tex, source lines 6327--6327. Cited by the retained statement or proof but not part of the collected spans; supplied only as a navigation pointer. The argument does not depend on its content.\par\medskip
\ifdefined\hypertarget\hypertarget{sourcepointer-7}{}\fi
\noindent\textit{Source reference.} Original lemma 0GH8: sources/stacks-artin/pinned-original-source.tex, source lines 5138--5153. Cited by the retained statement or proof but not part of the collected spans; supplied only as a navigation pointer. The argument does not depend on its content.\par\medskip
\ifdefined\hypertarget\hypertarget{sourcepointer-8}{}\fi
\noindent\textit{Source reference.} Original lemma 0A4Z: sources/stacks-artin/pinned-original-source.tex, source lines 6327--6327. Cited by the retained statement or proof but not part of the collected spans; supplied only as a navigation pointer. The argument does not depend on its content.\par\medskip
\begin{lemma}
\label{lemma-formal-object-effective}
In Situation \ref{situation-contractions} assume $X' \to S$ is separated.
Then every formal object for $F$ is effective.
\end{lemma}

\begin{proof}
A formal object $\xi = (R, \xi_n)$ of $F$ consists of a Noetherian
complete local $S$-algebra $R$ whose residue field is of finite type
over $S$, together with elements $\xi_n \in F(\Spec(R/\mathfrak m^n))$
for all $n$ such that $\xi_{n + 1}|_{\Spec(R/\mathfrak m^n)} = \xi_n$.
By the discussion in the proof of Lemma \ref{lemma-rs}
we see that either $\xi$ is a formal object of $U'$ or a formal
object of $W$. In the first case we see that $\xi$ is effective
by Lemma \ref{lemma-effective}. The second case is the interesting case.
Set $V = \Spec(R)$. We will construct an element
$(Z, u', \hat x) \in F(V)$ whose image in $F(\Spec(R/\mathfrak m^n))$
is $\xi_n$ for all $n \geq 1$.

\medskip\noindent
We may view the collection of elements $\xi_n$ as a morphism
$$
\xi : \text{Spf}(R) \longrightarrow W
$$
of locally Noetherian formal algebraic spaces over $S$. Observe that $\xi$
is {\it not} an adic morphism in general. To fix this, let $I \subset R$
be the ideal corresponding to the formal closed subspace
$$
\text{Spf}(R) \times_{\xi, W} W_{red} \subset \text{Spf}(R)
$$
Note that $I \subset \mathfrak m_R$. Set $Z = V(I) \subset V = \Spec(R)$.
Since $R$ is $\mathfrak m_R$-adically complete it is a fortiori
$I$-adically complete (Algebra, Lemma \ref{algebra-lemma-complete-by-sub}).
Moreover, we claim that for each $n \geq 1$ the morphism
$$
\xi|_{\text{Spf}(R/I^n)} :
\text{Spf}(R/I^n)
\longrightarrow
W
$$
actually comes from a morphism
$$
\xi'_n :  \Spec(R/I^n) \longrightarrow W
$$
Namely, this follows from writing $W = \colim W_n$ as in the
proof of Lemma \ref{lemma-rs}, noticing that $\xi|_{\text{Spf}(R/I^n)}$
maps into $W_n$, and applying Formal Spaces, Lemma
\ref{formal-spaces-lemma-map-into-algebraic-space}
to algebraize this to a morphism $\Spec(R/I^n) \to W_n$
as desired. Let us denote $\text{Spf}'(R) = V_{/Z}$ the formal spectrum
of $R$ endowed with the $I$-adic topology -- equivalently the formal
completion of $V$ along $Z$. Using the morphisms
$\xi'_n$ we obtain an adic morphism
$$
\hat x = (\xi'_n) : \text{Spf}'(R) \longrightarrow W
$$
of locally Noetherian formal algebraic spaces over $S$.
Consider the base change
$$
\text{Spf}'(R) \times_{\hat x, W, g} X'_{/T'} \longrightarrow \text{Spf}'(R)
$$
This is a formal modification by
Algebraization of Formal Spaces, Lemma
\ref{restricted-lemma-base-change-formal-modification}.
Hence by the main theorem on dilatations
(Algebraization of Formal Spaces, Theorem \ref{restricted-theorem-dilatations})
we obtain a proper morphism
$$
V' \longrightarrow V = \Spec(R)
$$
which is an isomorphism over $\Spec(R) \setminus V(I)$ and
whose completion recovers the formal modification above, in other words
$$
V' \times_{\Spec(R)} \Spec(R/I^n) =
\Spec(R/I^n) \times_{\xi'_n, W, g} X'_{/T'}
$$
This in particular tells us we have a compatible system of morphisms
$$
V' \times_{\Spec(R)} \Spec(R/I^n) \longrightarrow X' \times_S \Spec(R/I^n)
$$
Hence by Grothendieck's algebraization theorem (in the form of
More on Morphisms of Spaces, Lemma
\ref{spaces-more-morphisms-lemma-algebraize-morphism})
we obtain a morphism
$$
x' : V' \to X'
$$
over $S$ recovering the morphisms displayed above. Then finally
setting $u' : V \setminus Z \to X'$ the restriction of $x'$ to
$V \setminus Z \subset V'$ gives the third component of our
desired element $(Z, u', \hat x) \in F(V)$.
\end{proof}

\end{document}
